\section{A compact moduli space of bundles}\label{sec:moduli}

We first show that the rank condition in \eqref{eq:family} excludes
both singular sheaves and singular stable factors. The second
assertion is needed for the determinant polarization: the cycle
recorded by the Donaldson--Uhlenbeck morphism is the defect of the
slope graded object, not just the defect of the original sheaf.

\begin{proposition}\label{prop:no-bubbling}
Assume \eqref{eq:family}. Every slope-semistable torsion-free sheaf
of Mukai vector $2v_0$ is locally free. If it is strictly
slope-semistable, its slope Jordan--H\"older factors are two stable
bundles of vector $v_0$. Every slope-semistable sheaf of vector
$v_0$ is stable and locally free. For these two Mukai vectors,
slope and Gieseker semistability agree, as do stability and
S-equivalence.
\end{proposition}

\begin{proof}
We use two elementary properties of stable sheaves on a K3
surface. If $A$ is stable, simplicity and Serre duality give
\begin{equation}\label{eq:stable-bound}
 v(A)^2=\dim\Ext^1(A,A)-2\ge-2.
\end{equation}
If $A$ is slope-stable and torsion-free, $A^{**}$ is a stable
bundle, and
\[
 v(A^{**})=v(A)+(0,0,\operatorname{length}(A^{**}/A)).
\]
The reflexive hull of a slope-semistable sheaf is semistable.

Let $F$ have vector $2v_0$, and write
$\ell=\operatorname{length}(F^{**}/F)$. If $F^{**}$ is stable,
\eqref{eq:stable-bound} implies
\[
 4q-4r_0\ell\ge-2.
\]
Since $r_0\ge q+1$, this forces $\ell=0$.

Suppose instead that $F^{**}$ is strictly slope-semistable.
An equal-slope factor of rank $r'$ and first Chern class $a'H$
satisfies
\[
 r_0a'H^2=r'aH^2.
\]
Coprimality in \eqref{eq:family} forces $r_0\mid r'$.
There are therefore exactly two stable factors $A_1,A_2$, each
of rank $r_0$ and first Chern class $aH$. Put
\[
 v(A_i)=(r_0,aH,s_i),\qquad
 \ell_i=\operatorname{length}(A_i^{**}/A_i).
\]
Additivity gives $s_1+s_2=2s_0+\ell$. Applied to $A_i^{**}$,
\eqref{eq:stable-bound} gives
\begin{equation}\label{eq:factor-inequality}
 q-2r_0(s_i+\ell_i-s_0)\ge-2.
\end{equation}
Because $2r_0>q+2$, integrality implies
$s_i+\ell_i\le s_0$. Summing yields
$\ell+\ell_1+\ell_2\le0$. All three lengths vanish, and
$s_1=s_2=s_0$. Thus $F$ is a bundle and its stable factors are
bundles of vector $v_0$. Extensions of locally free sheaves are
locally free, consistently with this conclusion.

For vector $v_0$, the same coprimality rules out an equal-slope
proper factor. If $A$ is such a stable sheaf, its reflexive hull
has square $q-2r_0\ell$, and \eqref{eq:stable-bound} again forces
$\ell=0$.

Slope stability implies Gieseker stability, and Gieseker
semistability implies slope semistability. A strictly
slope-semistable $F$ of vector $2v_0$ has the two factors just
described. Their reduced Hilbert polynomials both equal that of
$F$, so $F$ is Gieseker-semistable and is not Gieseker-stable.
For vector $v_0$ there is no strictly slope-semistable case.
The Jordan--H\"older factors in the two theories consequently
coincide, proving the S-equivalence assertion.
\end{proof}

\subsection{Geometry of the moduli space}
The vector $v_0$ is primitive. On a surface of Picard rank one,
$H$ is general in the sense required for these moduli spaces
\cite[Definition~2.8 and Lemma~2.9(1)]{pr23}. The primitive
space $N$ is a nonempty smooth projective irreducible
holomorphic symplectic manifold of dimension
\begin{equation}\label{eq:n}
 n=q+2,
\end{equation}
deformation equivalent to $S^{[q/2+1]}$
\cite[Theorem~1.7 and the paragraph following Definition~1.6]{pr23}.
In particular, it is simply connected. Proposition~\ref{prop:no-bubbling}
identifies it with a compact moduli space of slope-stable bundles.

The results of Kaledin--Lehn--Sorger
\cite[Theorem~4.4, Proposition~6.1 and Theorems~A--B]{kls06}
give the following description. The variety $M$ is normal,
irreducible and projective, of dimension $4q+2$. Its regular
locus is precisely the nonempty stable locus. Every singular
point is represented by $G_1\oplus G_2$ with $G_i\in N$.
The direct-sum morphism identifies the reduced singular locus
with $\Sym^2N$. One can also see the last assertion from the
local quiver descriptions below: the semisimple nonsimple
representations are exactly the unordered pairs of simple
summands. Near a repeated summand, the restrictions of traces
of matrix words to two one-dimensional characters are the
polarized power sums. These generate the symmetric-square
invariant ring in characteristic zero, identifying the
reduced singular locus. Away from the diagonal the local
product description gives the same identification. Accordingly,
\begin{equation}\label{eq:strata}
 \Sing M=\Sigma=\Sym^2N,\qquad
 \Delta\cong N,\qquad
 \dim\Sigma=2q+4,\quad \dim\Delta=q+2.
\end{equation}
There are three smooth strata: $M^s$, $\Sigma\setminus\Delta$,
and $\Delta$.

The variety $M$ has symplectic singularities and is locally
factorial in the algebraic sense. Its singular locus has
codimension $2q-2\ge6$, so its singularities are terminal
\cite[Corollary~1]{namikawa01}. It has no proper symplectic
resolution \cite[Theorem~B]{kls06}. The Main Theorem of
\cite{namikawa06} applies to the projective, $\mathbb Q$-factorial,
terminal symplectic variety $M$ and makes every flat deformation
locally trivial. Since $\Sigma$ is nonempty, $M$ has no flat
smoothing. Algebraic local factoriality here does not assert
factoriality of completed local rings; see
\cite[Remark~6.3]{kls06}.

We shall also use the irreducible symplectic-variety theorem of
Perego--Rapagnetta \cite[Theorem~1.10]{pr23}. It gives canonical
singularities and
\begin{equation}\label{eq:forms-cohomology}
 H^0(M,\Omega_M^{[2]})=\CC\sigma_M,
 \qquad H^1(M,\OO_M)=0.
\end{equation}
For the second equality, the definition of an irreducible
symplectic variety gives no reflexive one-forms. Pullback to
a projective resolution, Hodge symmetry, and rationality of
canonical singularities give the stated cohomology vanishing;
compare \cite[Definition~1.2 and Remark~1.3]{pr23}.
Symplectic singularities are Gorenstein, with trivial canonical
bundle generated by the top exterior power of $\sigma_M$.

\subsection{The two extension cones}\label{subsec:quivers}
The formulas below use the harmonic Hermitian forms; the displayed
complex vector spaces are supplied with the corresponding flat
hyperk\"ahler structures. A quiver representation is called
simple if it has no nonzero proper subrepresentation. Let
$G_1\not\cong G_2$ belong to $N$.
Stability and Serre duality give
\[
 \Hom(G_1,G_2)=\Ext^2(G_1,G_2)=0,
 \qquad \dim\Ext^1(G_1,G_2)=q.
\]
Each self-extension space has dimension $n$. The effective
unitary stabilizer is a circle, acting with opposite weights
on the two cross-extension spaces. Hence
\begin{equation}\label{eq:offdiag-cone}
 \Cone_{G_1\oplus G_2}
 =\CC^{2n}\times\mathcal N_q,
 \qquad \mathcal N_q=T^*\CC^q\hkr_0\U(1).
\end{equation}
In coordinates $(x,y)\in\CC^q\oplus(\CC^q)^*$ the equations are
$y(x)=0$ and $|x|^2=|y|^2$. The map $(x,y)\mapsto x\otimes y$
identifies the complex quotient with the rank-at-most-one,
square-zero matrices in $\mathfrak{sl}_q$, namely the minimal
nilpotent orbit closure. Its complex dimension is $2q-2$.
There are simple balanced points: take nonzero perpendicular
vectors of equal norm for $x$ and the conjugate of $y$.

For $E=G\oplus G$, the extension space is
$\Ext^1(G,G)\otimes\mathfrak{gl}_2$. Splitting off scalar
matrices gives
\begin{equation}\label{eq:diagonal-cone}
 \Cone_{G\oplus G}=\CC^n\times\mathcal Q_q,
 \qquad
 \mathcal Q_q=(\CC^n\otimes\mathfrak{sl}_2)\hkr_0\SU(2),
 \qquad \dim\mathcal Q_q=3q.
\end{equation}
Here $n$ is even, and a symplectic basis identifies the extension
space with $n/2$ pairs of matrix loops $(A_j,B_j)$.
Simple balanced points exist: set all $B_j=0$ and choose two
noncommuting Hermitian traceless matrices among the $A_j$, with
the remaining matrices zero. The complex moment map
$\sum_j[A_j,B_j]$ and the real moment map
$\sum_j([A_j,A_j^*]+[B_j,B_j^*])$ vanish. The two Hermitian
matrices have no common invariant line, so this representation
is simple. Scaling makes these simple points arbitrarily small.
The scalar center acts trivially, so the use of $\SU(2)$ rather
than its effective quotient does not change the metric quotient.

The holomorphic germs and metric normalizations in
\eqref{eq:offdiag-cone}--\eqref{eq:diagonal-cone} will be identified
with those of $M$ in the next section. In particular, the metric
on the fixed scalar factor of \eqref{eq:diagonal-cone} is the
one induced by $a\mapsto a\otimes I_2$; it is twice the harmonic
metric on the deformation space of $G$.

\begin{example}\label{ex:quartic}
For a Picard-rank-one quartic K3 surface, $v_0=(5,H,0)$ satisfies
$q=4$, $r_0=q+1$ and $\gcd(5,4)=1$. The stable factors have
second Chern class $c_2=7$, and a sheaf of vector $2v_0$ has
$c_2=18$, since $s=r+c_1^2/2-c_2$ on a K3 surface. The
dimensions are as follows.
\begin{center}
\begin{adjustbox}{max width=\textwidth}
\begin{tabular}{lcc}
\toprule
Locus & Complex dimension & Complex normal-cone dimension\\
\midrule
$M^s$ & $18$ & $0$\\
$\Sigma\setminus\Delta$ & $12$ & $6$\\
$\Delta$ & $6$ & $12$\\
\bottomrule
\end{tabular}
\end{adjustbox}
\end{center}
Thus the family contains a compact example with both types of
singular points. The existence assertions are the moduli
theorems above; the numerical equalities alone do not establish
existence.
\end{example}
